\section{形式化验证：有限变换半群上的三生成元残差闭包}\label{app:formal-verification-lean}

本附录给出 \S\ref{subsec:pom-three-generator-residual-closure} 中三生成元残差闭包理论在有限变换半群 $T_n$（$\mathrm{Fin}\,n\to\mathrm{Fin}\,n$ 的全体函数在复合下构成的半群）上的一个完全可执行的具体化。
全部代码以 Lean~4（v4.27.0）实现，源文件位于本论文目录下 \texttt{formal/} 子目录，可由 \texttt{lake build} 复现编译。

\subsection{系统概述}

\paragraph{核心设定}
固定状态空间大小 $n$，考虑全变换半群 $T_n = \mathrm{End}(\mathrm{Fin}\,n)$，其中每个元素 $f\in T_n$ 可表示为长度为 $n$ 的数组 $[f(0),f(1),\dots,f(n{-}1)]$。
定义：
\begin{itemize}
  \item \textbf{秩}（rank）：$\mathrm{rank}(f):=|\mathrm{Im}(f)|$，即像集大小。
  \item \textbf{亏值}（defect）：$\mathrm{def}(f):=n-\mathrm{rank}(f)$，即信息损失量。
\end{itemize}

\paragraph{自动生成器搜索}
系统通过穷举搜索（\texttt{AutoSeed.findMinGenerators}）寻找能生成全部 $|T_n|=n^n$ 个元素的最小生成器集合。
对 $n\ge 3$ 的所有已测案例，最小生成器数量稳定为~$3$。

\paragraph{与论文三生成元的对应}
令自动发现的三个生成器为 $g_0,g_1,g_2$（按字典序排序），则对 $n=3$ 和 $n=4$，其结构如下：

\begin{center}
\begin{tabular}{c|ccc|c}
\toprule
$n$ & $g_0$ & $g_1$ & $g_2$ & 结构 \\
\midrule
3 & $[2,0,1]$ & $[2,1,0]$ & $[2,2,1]$ & 2个置换 $+$ 1个合并 \\
4 & $[3,1,2,0]$ & $[3,2,0,1]$ & $[3,3,2,1]$ & 2个置换 $+$ 1个合并 \\
\bottomrule
\end{tabular}
\end{center}

\noindent
稳定结构为：$g_0,g_1$ 均为满秩（$\mathrm{rank}=n$，置换），$g_2$ 的秩为 $n{-}1$（恰好将两个状态合并为一个）。
这与命题~\ref{prop:pom-three-gen-interface-normal-form} 的三生成元片段在结构上对应：

\begin{center}
\begin{tabular}{c|c|c}
\toprule
论文生成元 & HyperKernel 生成元 & 性质 \\
\midrule
$\mathrm{PROJ}_u,\;\mathrm{LIFT}_G$（可逆/可交换） & $g_0,\;g_1$（置换） & 不损失信息 \\
$E_m$（带残差） & $g_2$（合并） & 导致秩下降（信息损失）\\
\bottomrule
\end{tabular}
\end{center}

\subsection{异常签名与零异常比例}

\paragraph{定义}
对闭包字典中的每个函数 $f$，其最短生成词为 $w=g_{i_1}\cdots g_{i_\ell}$。
定义 $f$ 的\textbf{奇异计数}（singular count）为词 $w$ 中 $g_2$ 出现的次数：
$$
\mathrm{sc}(f):=\#\{j : i_j = 2\}.
$$
定义\textbf{excess}（对应论文的异常签名 $\mathsf{Anom}$）为：
$$
\mathrm{excess}(f):=\mathrm{sc}(f)-\mathrm{def}(f).
$$

\begin{proposition}[奇异计数下界——穷举验证]\label{prop:formal-singular-lower-bound}
对 $n=3$（$|T_3|=27$）和 $n=4$（$|T_4|=256$），穷举验证每一个 $f\in T_n$ 均满足
$$
\mathrm{sc}(f)\ge \mathrm{def}(f),
\qquad\text{即}\qquad
\mathrm{excess}(f)\ge 0.
$$
\end{proposition}

\begin{proof}
由 Lean~4 程序 \texttt{HyperKernel.Analysis.verifyDefectHypothesis} 对 $T_n$ 的全部 $n^n$ 个元素逐一核验得出。
\end{proof}

\paragraph{零异常函数的比例}
称 $f$ 为\textbf{零异常函数}（zero-anomaly），若 $\mathrm{excess}(f)=0$，即达到了信息论下界。
穷举统计结果如下：

\begin{center}
\begin{tabular}{c|c|c|c|c}
\toprule
$n$ & $|T_n|$ & 匹配（$\mathrm{excess}=0$） & 反例（$\mathrm{excess}>0$） & 零异常比例 \\
\midrule
3 & 27 & 24 & 3 & 88.9\% \\
4 & 256 & 184 & 72 & 71.9\% \\
\bottomrule
\end{tabular}
\end{center}

\noindent
这一比例在 \S\ref{subsec:pom-three-generator-residual-closure} 的框架下可解读为：在有限变换半群上，\textbf{绝大多数函数的最短实现恰好使用了信息论所要求的最少不可逆操作次数}，仅有少部分函数存在"多余的"不可逆步骤（正异常）。

\subsection{秩分布与复杂度反比律}

\paragraph{秩分布}
函数按秩的分布展示了 $T_n$ 的组合结构：

\begin{center}
\begin{tabular}{c|cccc}
\toprule
$n=4$ & $\mathrm{rank}=1$ & $\mathrm{rank}=2$ & $\mathrm{rank}=3$ & $\mathrm{rank}=4$ \\
\midrule
函数数 & 4 & 84 & 144 & 24 \\
平均词长 & 6.00 & 5.87 & 5.13 & 3.58 \\
\bottomrule
\end{tabular}
\end{center}

\begin{center}
\begin{tabular}{c|ccc}
\toprule
$n=3$ & $\mathrm{rank}=1$ & $\mathrm{rank}=2$ & $\mathrm{rank}=3$ \\
\midrule
函数数 & 3 & 18 & 6 \\
平均词长 & 3.67 & 3.17 & 1.33 \\
\bottomrule
\end{tabular}
\end{center}

\begin{proposition}[秩–复杂度反比律——穷举验证]\label{prop:formal-rank-complexity}
对 $n=3$ 和 $n=4$，函数的平均最短词长关于秩严格递减：
$$
\mathrm{rank}(f) < \mathrm{rank}(f') \implies \EE[|w| \mid \mathrm{rank}=\mathrm{rank}(f)] > \EE[|w| \mid \mathrm{rank}=\mathrm{rank}(f')].
$$
即秩越低（信息损失越多）的函数需要越长的生成词才能实现。
\end{proposition}

\subsection{反例的结构分析}

\paragraph{$n=3$ 的反例}
$n=3$ 仅有 $3$ 个反例，均具有 $\mathrm{rank}=2$、$\mathrm{def}=1$、$\mathrm{sc}=2$：

\begin{center}
\begin{tabular}{c|c|c|c|c}
\toprule
函数 & 秩 & 亏值 & 奇异计数 & 最短词 \\
\midrule
$[1,2,1]$ & 2 & 1 & 2 & $g_0 g_0 g_2 g_2$ \\
$[2,1,1]$ & 2 & 1 & 2 & $g_0 g_2 g_2$ \\
$[1,1,2]$ & 2 & 1 & 2 & $g_2 g_2$ \\
\bottomrule
\end{tabular}
\end{center}

\noindent
共同特征：所有反例的最短词中包含两次连续的 $g_2$，但亏值仅为~$1$——这意味着第二次 $g_2$ 并未进一步损失信息（秩未再下降），但在词法上却不可省去。
这对应于 \S\ref{subsec:pom-anom-cohomology} 中异常签名的"不可消去"性质：即便交换失败在语义上没有额外损失信息，其词法残差仍须被忠实记录。

\paragraph{$n=4$ 的反例分类}
$n=4$ 共有 $72$ 个反例（占 $28.1\%$）。
其中展示前 $10$ 个的典型结构：

\begin{center}
\begin{tabular}{c|c|c|c|c}
\toprule
函数 & 秩 & 亏值 & 奇异计数 & 词长 \\
\midrule
$[3,0,0,3]$ & 2 & 2 & 3 & 9 \\
$[3,1,1,3]$ & 2 & 2 & 3 & 8 \\
$[0,3,3,0]$ & 2 & 2 & 3 & 8 \\
$[1,2,2,1]$ & 2 & 2 & 4 & 8 \\
$[1,0,0,1]$ & 2 & 2 & 3 & 8 \\
$[0,2,2,0]$ & 2 & 2 & 3 & 8 \\
$[0,3,0,3]$ & 2 & 2 & 3 & 8 \\
$[3,3,0,0]$ & 2 & 2 & 3 & 7 \\
$[2,1,1,2]$ & 2 & 2 & 4 & 8 \\
$[0,1,1,0]$ & 2 & 2 & 3 & 8 \\
\bottomrule
\end{tabular}
\end{center}

\noindent
\textbf{观察}：所有反例均具有 $\mathrm{rank}=2$（$\mathrm{def}=2$），且 $\mathrm{sc}\ge 3$（excess $\ge 1$）。
这表明正异常集中出现在"中等亏值"区域；满秩（置换）和最低秩（常值函数）的函数均为零异常。

\subsection{Cayley 图直径与统计摘要}

完整统计摘要：

\begin{center}
\begin{tabular}{c|cccccccc}
\toprule
$n$ & $|T_n|$ & 直径 & 平均词长 & 平均秩 & 平均亏值 & 平均 $g_0$ & 平均 $g_1$ & 平均 $g_2$ \\
\midrule
3 & 27 & 5 & 2.81 & 2.11 & 0.89 & 1.44 & 0.37 & 1.00 \\
4 & 256 & 9 & 5.24 & 2.73 & 1.27 & 1.41 & 2.26 & 1.57 \\
\bottomrule
\end{tabular}
\end{center}

\subsection{与论文理论的对齐总结}\label{subsec:formal-alignment-summary}

以下表格总结 HyperKernel 形式化验证结果与本文 \S\ref{subsec:pom-three-generator-residual-closure} 理论框架的对齐关系：

\begin{center}
\begin{tabular}{p{5cm}|p{5cm}|p{4cm}}
\toprule
论文概念 & HyperKernel 实现 & 验证状态 \\
\midrule
三生成元稳定性 & $n\ge 3$: 最小生成器数${}=3$ & 穷举验证 ($n=3,4$) \\
$2$个可交换 $+$ $1$个带残差 & $2$个置换 $+$ $1$个合并 & 穷举验证 \\
异常签名 $\mathsf{Anom}_G$ & $\mathrm{excess}(f)=\mathrm{sc}(f)-\mathrm{def}(f)$ & 穷举计算 \\
零异常占多数 & $88.9\%$ ($n{=}3$), $71.9\%$ ($n{=}4$) & 穷举统计 \\
下界 $\mathrm{sc}\ge\mathrm{def}$ & 全部 $n^n$ 个函数满足 & 穷举验证 \\
秩–复杂度反比律 & 平均词长随秩严格递减 & 穷举统计 \\
\bottomrule
\end{tabular}
\end{center}

\subsection{复现指南}

全部源代码位于 \texttt{formal/} 目录，需要 Lean~4 v4.27.0 环境。

\begin{verbatim}
cd formal/
lake build              # 编译全部模块
lake exe hyperkernel    # 运行基础闭包计算
lake exe analyze        # 运行深度结构分析
\end{verbatim}

\noindent
修改 \texttt{formal/HyperKernel/Spec.lean} 中的 \texttt{n} 值可切换状态空间大小。
$n=3$ 在数秒内完成，$n=4$ 约需 $1{-}2$ 分钟。

\subsection{重写系统分析（论文定义 \ref{def:pom-three-gen-rewrite-rules} 的有限半群具体化）}

\paragraph{方法}
论文定义~\ref{def:pom-three-gen-rewrite-rules} 给出的重写规则族要求将投影词排序为正规序 $\mathrm{LIFT}\prec\mathrm{PROJ}\prec E$。
在有限半群中，我们对每个函数的 BFS 最短词执行"冒泡排序式"的重写：逐步将相邻的逆序生成元对 $(g_i, g_j)$（$i>j$）交换为 $(g_j, g_i)$，并记录每次交换是否改变了词所计算的函数。

\begin{definition}[逆序度量（命题 \ref{prop:pom-three-gen-termination} 的具体化）]\label{def:formal-inversion-measure}
对词 $w = g_{i_1}\cdots g_{i_\ell}$，定义逆序度量
$$
\mu(w) := \#\{(a,b) : a < b,\; i_a > i_b\}.
$$
每次冒泡交换严格降低 $\mu$，因此重写过程在 $O(\ell^2)$ 步内必然终止。
\end{definition}

\paragraph{结果（$n=4$）}

\begin{center}
\begin{tabular}{l|r}
\toprule
指标 & 值 \\
\midrule
函数总数 & 256 \\
已在正规形（无需重写）& 19 \\
重写后函数不变（可交换） & 19 \\
重写后函数改变（不可交换） & 237 \\
平均重写步数 & 4.21 \\
平均交换失败次数 & 4.21 \\
最大交换失败次数 & 12 \\
\bottomrule
\end{tabular}
\end{center}

\begin{remark}[非交换性是普遍现象]\label{rem:formal-noncommutativity}
在 $256$ 个函数中仅 $19$ 个（$7.4\%$）的最短词在冒泡排序下保持函数不变。
这与论文 \S\ref{subsec:pom-three-generator-residual-closure} 的框架一致：在有限半群 $T_n$ 中，生成元之间的交换在绝大多数情况下产生非零残差——\textbf{交换失败是常态，而非例外}。
每一次使函数改变的交换步对应论文中的一次"异常累积"（定义~\ref{def:pom-anom-accumulation}）。
\end{remark}

\subsection{正规形可达性分析（论文命题 \ref{prop:pom-three-gen-interface-normal-form} 的有限半群具体化）}

\paragraph{方法}
论文命题~\ref{prop:pom-three-gen-interface-normal-form} 断言：任意三生成元投影词都可重写为正规形 $\mathrm{LIFT}_G\circ\mathrm{PROJ}_u\circ E_m$。
在有限半群中，正规形对应"排序词"$g_0^a g_1^b g_2^c$（所有 $g_0$ 在前，$g_1$ 居中，$g_2$ 在后）。
我们穷举长度 $\le 11$（直径 $+2$）的所有排序词，检查 $T_4$ 的每个函数是否可由某个排序词生成。

\paragraph{结果（$n=4$）}

\begin{center}
\begin{tabular}{l|r}
\toprule
指标 & 值 \\
\midrule
函数总数 & 256 \\
可达（存在排序词生成） & 24 \\
不可达 & 232 \\
可达率 & 9.4\% \\
\bottomrule
\end{tabular}
\end{center}

\paragraph{按秩的可达率}

\begin{center}
\begin{tabular}{c|c|c}
\toprule
秩 & 可达/总数 & 可达率 \\
\midrule
$\mathrm{rank}=1$ & 0/4 & 0\% \\
$\mathrm{rank}=2$ & 0/84 & 0\% \\
$\mathrm{rank}=3$ & 16/144 & 11.1\% \\
$\mathrm{rank}=4$ & 8/24 & 33.3\% \\
\bottomrule
\end{tabular}
\end{center}

\begin{proposition}[有限半群中的正规形不可达性]\label{prop:formal-nf-unreachable}
在 $T_4$ 中，$90.6\%$ 的函数不可由任何排序词 $g_0^a g_1^b g_2^c$ 生成。
特别地，所有 $\mathrm{rank}\le 2$ 的函数（$88$ 个）均不可达。
\end{proposition}

\begin{remark}[与论文理论的关键差异：语义正规形 vs.\ 词法正规形]\label{rem:formal-semantic-vs-syntactic}
论文命题~\ref{prop:pom-three-gen-interface-normal-form} 的正规形定理建立在\textbf{语义层}：生成元 $\mathrm{PROJ}_u$、$\mathrm{LIFT}_G$、$E_m$ 携带连续参数（温度 $u$、群 $G$、分辨率 $m$），因此即使词被排序后参数改变，仍可在连续参数空间中寻找新的参数使得排序词表达同一语义。
相比之下，有限半群 $T_n$ 中的生成元是\textbf{固定的离散对象}，不携带连续参数——交换两个生成元在词中的位置一般会改变所计算的函数。

这一差异恰好量化了论文中"异常签名"的信息内容：在连续参数设定下，交换失败可被参数调整吸收（推论~\ref{cor:pom-zero-anom-pullback-exactness}）；在离散设定下，交换失败无法被吸收，因此表现为\textbf{不可达性}——这是"离散异常"的最极端形式。
\end{remark}

\begin{corollary}[离散异常的层级结构]\label{cor:formal-discrete-anomaly-hierarchy}
将正规形不可达性视为"极端异常"，有限半群上的异常呈现如下层级：
\begin{enumerate}
  \item \textbf{零异常函数}：$\mathrm{excess}=0$ 且可由排序词生成——$18$ 个（$7.0\%$）。
  \item \textbf{弱异常函数}：$\mathrm{excess}=0$ 但不可由排序词生成——$166$ 个（$64.8\%$）。
  \item \textbf{强异常函数}：$\mathrm{excess}>0$（必然不可由排序词生成）——$72$ 个（$28.1\%$）。
\end{enumerate}
此层级结构表明：即使奇异计数达到了信息论下界（$\mathrm{sc}=\mathrm{def}$），生成元的非交换性仍然阻止了大部分函数被正规形词表达。
\end{corollary}

\subsection{源文件清单}

\begin{center}
\begin{tabular}{l|p{9cm}}
\toprule
文件 & 功能 \\
\midrule
\texttt{Op.lean} & 有限态操作 $\mathrm{Op}\,n$ 的定义（数组表示、复合、相等判定）\\
\texttt{Spec.lean} & 参数配置（状态空间大小 $n$、搜索上界 \texttt{maxSeedSize}）\\
\texttt{Enum.lean} & $T_n$ 的完全枚举 \\
\texttt{Closure.lean} & BFS 闭包计算与最短词字典 \\
\texttt{AutoSeed.lean} & 最小生成器集合的自动搜索 \\
\texttt{Pretty.lean} & 格式化输出与排序工具 \\
\texttt{Run.lean} & 主执行逻辑（基础版） \\
\texttt{Analysis.lean} & 深度结构分析（秩、亏值、奇异计数、假设验证）\\
\texttt{Rewrite.lean} & 重写系统（冒泡排序式正规化、逆序度量、交换失败检测）\\
\texttt{NormalForm.lean} & 正规形分析（排序词可达性、长度开销计算）\\
\texttt{AnalyzeRun.lean} & 分析执行逻辑（统计、分布、反例展示）\\
\bottomrule
\end{tabular}
\end{center}
